\scanpage{14}
1) Soit une extension de groupes topologiques
\[
1\to N\to G\to H\to1 \tag*{(E)}
\]
et un hom de groupes topologiques
\[
 H'\xrightarrow{u}H
\]
d'où par image inverse une extension \inserted{$G'$} de $H'$ par $N$ et un hom d'extensions
\diagram{\begin{tikzcd}[ampersand replacement=\&,column sep=1.25em,row sep=1.5em]
(E) \& 1 \arrow[r] \& N \arrow[r] \& G \arrow[r] \& H \arrow[r] \& 1\\
(E') \& 1 \arrow[r] \& N^{\deleted{\unsure{\prime}}} \arrow[u,"\mathrm{id}"'] \arrow[r] \& G' \arrow[u] \arrow[r] \& H' \arrow[u,"u"'] \arrow[r] \& 1
\end{tikzcd}}
Passant aux $\pi_0$, on en déduit un hom de suites exactes \deleted{(sans $1$ à gauche}
\diagram{\begin{tikzcd}[ampersand replacement=\&,column sep=1.25em,row sep=1.5em]
(S) \& \pi_1 H \arrow[r] \& \pi_0N \arrow[r] \& \pi_0G \arrow[r] \& \pi_0H \arrow[r] \& 1\\
(S') \& \pi_1 H' \arrow[u] \arrow[r] \& \pi_0N \arrow[u,equal] \arrow[r] \& \pi_0G' \arrow[u] \arrow[r] \& \pi_0H' \arrow[u] \arrow[r] \& 1
\end{tikzcd}}
dont le carré de droite donne lieu à un hom
\[
\underset{\substack{\Vert\\\pi_0(G\times_H H')}}{\pi_0G'}\xrightarrow{\unsure{\varepsilon|c}}
\pi_0G\times_{\pi_0H}\pi_0H'.
\]

Il est facile de voir sur des exemples que cet hom (qui est tjrs surjectif) n'est pas néc. bijectif. Ainsi, la seule connaissance de \inserted{(disons)} la suite $(S)$ et des $\pi_i(H')\to\pi_i(H)$ ne permet pas, en général, de \deleted{\unsure{calculer}} déterminer $\pi_0(G')$. Notre propos est de préciser des données algébriques simples, \unsure{rel$^{\rm t}$es|relatives} à l'extension $(E)$, qui permettent de \deleted{calculer} \inserted{déterminer} \inserted{(mod isom)} $\pi_0(G')$ connaissant $\pi_0\unsure{(H)|H}\to\pi_0H'$ et $\pi_1H\to\pi_1H'$ \unsure{induits|induit} par $u$.
% archive sheet number: 13
